\documentclass[10pt]{article}
\usepackage[a4paper,margin=0.86in]{geometry}
\renewcommand{\rmdefault}{cmr}
\usepackage{amsmath,amssymb,amsthm,mathtools,booktabs,array,enumitem,xcolor,hyperref}
\hypersetup{colorlinks=true,linkcolor=blue!55!black,citecolor=blue!55!black,urlcolor=blue!55!black}
\setlength{\emergencystretch}{3em}
\setlength{\parindent}{0pt}
\setlength{\parskip}{0.52em}
\setlist{leftmargin=1.55em,itemsep=0.22em,topsep=0.24em}
\sloppy

\newtheorem{theorem}{Theorem}[section]
\newtheorem{proposition}[theorem]{Proposition}
\newtheorem{lemma}[theorem]{Lemma}
\newtheorem{corollary}[theorem]{Corollary}
\theoremstyle{definition}
\newtheorem{definition}[theorem]{Definition}
\newtheorem{convention}[theorem]{Convention}
\newtheorem{remark}[theorem]{Remark}

\newcommand{\Lam}{\Lambda}
\newcommand{\Sh}{\mathsf{Sh}}
\newcommand{\Mod}{\mathsf{Mod}}
\newcommand{\Cpx}{\mathsf{C}}
\newcommand{\Kcat}{\mathsf{K}}
\newcommand{\Cons}{\mathsf{Cons}}
\newcommand{\et}{\mathrm{et}}
\newcommand{\proet}{\mathrm{proet}}
\newcommand{\Dbc}{D^b_c}
\newcommand{\Dctf}{D^b_{\mathrm{ctf}}}
\newcommand{\DbcAt}[1]{D^b_{c,#1}}
\newcommand{\DctfAt}[1]{D^b_{\mathrm{ctf},#1}}
\newcommand{\Tor}{\operatorname{Tor}}
\newcommand{\Hom}{\operatorname{Hom}}
\newcommand{\Cone}{\operatorname{Cone}}
\newcommand{\id}{\operatorname{id}}
\newcommand{\dep}[1]{\par\smallskip\noindent\textit{Direct dependencies: #1}}

\begin{document}
\begin{center}
{\large\bfseries Finite Torsion Sheaves on the Etale and Pro-Etale Sites\par}
\medskip
{\normalsize Sites, constructibility, and derived comparison\par}
\end{center}

\begin{abstract}
This note isolates the finite-torsion foundation required before any adic
completion is formed.  It gives the sites, coefficient sheaves, constructible
hearts, the finite etale--pro-etale comparison, Verdier-localized derived
categories, the standard $t$-structure, derived tensor products, and scheme
pullback.  The comparison theorem is decomposed into its sheaf-theoretic,
cohomological, and derived constituents.  All derived categories below are
ordinary triangulated categories.  No direct-image theorem for a scheme
morphism, compact-support construction, duality theorem, or trace formula is
included here.
\end{abstract}

\tableofcontents

\section{Coefficient and size data}

\begin{convention}[Universes and selected sites]
Fix Grothendieck universes $\mathcal U\in\mathcal V$.  Schemes, rings, modules,
and test schemes are taken in $\mathcal U$.  Each site is replaced by a
$\mathcal U$-small skeleton that contains the final object and is closed, up to
the chosen representatives, under finite fiber products.  Categories of
sheaves and their localizations are regarded as $\mathcal V$-categories.
The chosen skeleton is required to be a basis: every admissible test scheme is
covered by selected objects, and every covering of a selected object admits a
cofinal selected refinement.  Changing to another sufficiently large selected
basis produces an equivalent sheaf category.
\end{convention}

\begin{definition}[Finite torsion coefficient datum]
A finite torsion coefficient datum on a scheme $X$ consists of a finite
commutative ring $\Lam$ and an integer $N>0$ such that $N\Lam=0$.  The datum is
\emph{invertible on $X$} if $N$ is a unit in every local ring
$\mathcal O_{X,x}$, equivalently if $X\to\operatorname{Spec}\mathbf Z$ factors
through $\operatorname{Spec}\mathbf Z[1/N]$.

Finiteness of $\Lam$ is used throughout this note.  Invertibility is not needed
to define the categories or the etale--pro-etale comparison; it becomes a
hypothesis for later cohomological finiteness and base-change theorems.  The
principal example is $\Lam=\mathbf Z/\ell^n\mathbf Z$ with $\ell$ invertible on
$X$.
\end{definition}

\begin{lemma}[Noetherianity of the coefficient ring]
Every ideal of $\Lam$ and every submodule of a finite $\Lam$-module is finite
and finitely generated.  In particular, the finite ring $\Lam$ is noetherian.
\dep{the finite torsion coefficient datum.}
\end{lemma}

\section{The small etale and pro-etale sites}

\begin{definition}[Weakly etale morphism]
A morphism of schemes $u:U\to X$ is \emph{weakly etale} if $u$ is flat and its
diagonal
\[
\Delta_u:U\longrightarrow U\times_XU
\]
is flat.  This condition is stable under composition and base change.
\end{definition}

\begin{definition}[Selected small etale site]
The site $X_{\et}$ has selected etale morphisms $U\to X$ as objects and
$X$-morphisms as arrows.  A family $\{U_i\to U\}_{i\in I}$ is a covering if
each $U_i\to U$ is etale and the images jointly cover $|U|$.
\end{definition}

\begin{definition}[Selected small pro-etale site]
The site $X_{\proet}$ has selected weakly etale morphisms $U\to X$ as objects
and $X$-morphisms as arrows.  A family $\{U_i\to U\}_{i\in I}$ is a covering
if its arrows are weakly etale and, for every affine open $V\subseteq U$, there
are finitely many affine opens $V_j\subseteq U_{i(j)}$ whose images cover $V$.
In particular, every such covering is fpqc and jointly surjective.
\end{definition}

\begin{proposition}[Site axioms and subcanonicity]
For $\tau\in\{\et,\proet\}$, the selected category $X_\tau$ has a final object
and finite fiber products.  Coverings are stable under pullback and composition,
and every representable presheaf is a sheaf.  Hence $X_\tau$ is a small
subcanonical site.
\dep{the universe convention and Definitions 2.1--2.3; etale descent and fpqc
descent for morphisms.}
\end{proposition}

\begin{definition}[Presheaves, sheaves, and sheafification]
Let $(\mathcal C,J)$ be either selected site.  A presheaf of sets is a
contravariant functor $F:\mathcal C^{\mathrm{op}}\to\mathsf{Set}$.  It is a
sheaf if, for every covering $\{U_i\to U\}$, the diagram
\[
F(U)\longrightarrow\prod_iF(U_i)
\rightrightarrows\prod_{i,j}F(U_i\times_UU_j)
\]
is an equalizer.  Sheaves of abelian groups, rings, and modules are defined by
the same condition on their underlying set-valued presheaves.  Sheafification
$a_\tau P$ is characterized by a universal map $P\to a_\tau P$ from a
presheaf to a sheaf.
\end{definition}

\begin{definition}[Finite constant coefficient sheaf]
Let $c_X\Lam$ be the presheaf on $X_\tau$ that is constantly $\Lam$ with
identity restriction maps.  Define
\[
\underline\Lam_{X,\tau}:=a_\tau(c_X\Lam).
\]
Equivalently, $\underline\Lam_{X,\tau}(U)$ consists of locally constant maps
from $U$ to the finite discrete ring $\Lam$; on a pro-etale object these are
the continuous locally constant maps $|U|\to\Lam$.
\end{definition}

\begin{definition}[Module-sheaf category]
An $\underline\Lam_{X,\tau}$-module sheaf is an abelian sheaf $M$ together
with a unital associative action
\[
\underline\Lam_{X,\tau}\otimes_{\mathbf Z}M\longrightarrow M.
\]
Morphisms commute with the action.  Write
\[
\mathsf A_\tau(X,\Lam)
:=\Mod(X_\tau,\underline\Lam_{X,\tau}).
\]
\end{definition}

\begin{theorem}[Ambient Grothendieck abelian category]
For $\tau\in\{\et,\proet\}$, $\mathsf A_\tau(X,\Lam)$ is a Grothendieck
abelian category.  Kernels are computed objectwise, cokernels are the
module-sheafifications of objectwise cokernels, filtered colimits are exact,
and the category has enough injectives and a generator.
\dep{Proposition 2.4, sheafification on a small site, and Definition 2.7.}
\end{theorem}

\begin{definition}[Geometric stalk]
A geometric point is a morphism $\bar x:\operatorname{Spec}\Omega\to X$ with
$\Omega$ separably closed.  For an etale sheaf $F$, its stalk is the filtered
colimit over pointed etale neighborhoods of $\bar x$.  For a pro-etale sheaf,
one may use pointed pro-etale neighborhoods; the resulting fiber functor is
denoted in both cases by $F\mapsto F_{\bar x}$.
\end{definition}

\section{Finite locally constant and constructible sheaves}

\begin{definition}[Finite locally constant module sheaf]
An object $F\in\mathsf A_\tau(X,\Lam)$ is \emph{finite locally constant} if
there is a covering $\{U_a\to X\}$ in $X_\tau$ and finite $\Lam$-modules
$M_a$ such that
\[
F|_{U_a}\simeq\underline{M_a}_{U_a,\tau}
\]
for every $a$.  Since $\Lam$ is noetherian, ``finite'' may equivalently be
replaced by ``finitely presented'' here.
\end{definition}

\begin{lemma}[Common trivializing cover]
Let $I$ be a finite category and let $F:I\to\mathsf A_\tau(X,\Lam)$ be a
diagram whose values are finite locally constant.  There is one covering of
$X$ on which every $F(i)$ is constant and every arrow of the diagram is
induced by a homomorphism of finite $\Lam$-modules.
\dep{Definition 3.1, finite fiber products in Proposition 2.4, and finiteness
of the modules.}
\end{lemma}

\begin{proposition}[Closure of finite local systems]
For a morphism $u:F\to G$ of finite locally constant
$\underline\Lam_{X,\tau}$-modules, $\ker u$ and $\operatorname{coker}u$ are
finite locally constant.  If
\[
0\longrightarrow F\longrightarrow E\longrightarrow G\longrightarrow0
\]
is exact and $F,G$ are finite locally constant, then $E$ is finite locally
constant.  Finite direct sums, tensor products, and direct summands of finite
locally constant sheaves are finite locally constant.
\dep{Lemma 3.2, Lemma 1.3, and exactness of inverse image on a site.}
\end{proposition}

\begin{definition}[Constructible subset and stratification]
A subset of a scheme is constructible if it is a finite union of locally
closed quasi-compact subsets.  A finite constructible stratification of a
scheme $V$ is a finite disjoint decomposition
\[
V=\coprod_{\alpha\in A}V_\alpha
\]
into reduced locally closed subschemes whose underlying subsets are
constructible.
\end{definition}

\begin{definition}[Constructible finite-torsion sheaf]
An object $F\in\mathsf A_\tau(X,\Lam)$ is \emph{constructible} if, for every
affine open $V\subseteq X$, there is a finite constructible stratification
$V=\coprod_\alpha V_\alpha$ such that $F|_{V_\alpha}$ is finite locally
constant.  Let
\[
\Cons_\tau(X,\Lam)\subseteq\mathsf A_\tau(X,\Lam)
\]
be the strictly full subcategory of constructible sheaves.
\end{definition}

\begin{proposition}[Global finite stratifications]
If $X$ is quasi-compact and quasi-separated, the affine-local condition in
Definition 3.5 is equivalent to the existence of one finite constructible
stratification of $X$ adapted to $F$.  Any finite collection of such
stratifications has a finite common refinement.
\dep{Definitions 3.4 and 3.5; stability of constructible locally closed subsets
under finite intersections.}
\end{proposition}

\begin{theorem}[Constructible kernels, cokernels, and extensions]
Let $u:F\to G$ be a morphism in $\mathsf A_\tau(X,\Lam)$ with $F,G$
constructible.  Then $\ker u$ and $\operatorname{coker}u$ are constructible.
If $0\to F\to E\to G\to0$ is exact with $F,G$ constructible, then $E$ is
constructible.
\dep{Proposition 3.3, Definition 3.5, and the common-refinement assertion of
Proposition 3.6, applied on every affine open.}
\end{theorem}

\begin{corollary}[The constructible heart]
The category $\Cons_\tau(X,\Lam)$ is abelian.  Its inclusion into
$\mathsf A_\tau(X,\Lam)$ is exact and extension closed.  Ambient finite limits
and finite colimits of constructible sheaves remain constructible.
\dep{Theorems 2.8 and 3.7.}
\end{corollary}

\begin{proposition}[Tensor closure in the heart]
The sheaf tensor product defines a bifunctor
\[
-\otimes_{\underline\Lam_{X,\tau}}-:
\Cons_\tau(X,\Lam)\times\Cons_\tau(X,\Lam)
\longrightarrow\Cons_\tau(X,\Lam).
\]
It is right exact in each variable and carries finite locally constant
sheaves to finite locally constant sheaves.
\dep{Lemma 3.2, Proposition 3.3, and Proposition 3.6.}
\end{proposition}

\section{Scheme pullback before deriving}

\begin{definition}[Morphism of sites and same-coefficient pullback]
Let $f:Y\to X$ be a morphism of schemes.  For
$\tau\in\{\et,\proet\}$, the induced geometric morphism of topoi is denoted
\[
f_\tau:\Sh(Y_\tau)\longrightarrow\Sh(X_\tau).
\]
The coefficient morphism
$f_\tau^{-1}\underline\Lam_{X,\tau}\to\underline\Lam_{Y,\tau}$ is an
isomorphism.  Define
\[
f_\tau^*M:=
\underline\Lam_{Y,\tau}\otimes_{f_\tau^{-1}\underline\Lam_{X,\tau}}
f_\tau^{-1}M,
\qquad
f_\tau^*: \mathsf A_\tau(X,\Lam)\to\mathsf A_\tau(Y,\Lam).
\]
Under the coefficient isomorphism this is simply inverse image.
\end{definition}

\begin{proposition}[Exactness and monoidality of pullback]
For every morphism $f:Y\to X$, the functor $f_\tau^*$ is exact, preserves all
colimits and finite limits, and is strong symmetric monoidal.  Thus there are
natural isomorphisms
\[
f_\tau^*(M\otimes N)\simeq f_\tau^*M\otimes f_\tau^*N,
\qquad f_\tau^*\underline\Lam_{X,\tau}\simeq
\underline\Lam_{Y,\tau}.
\]
\dep{Definition 4.1 and the exact inverse-image functor of a geometric
morphism of topoi.}
\end{proposition}

\begin{proposition}[Constructibility under pullback]
Assume $f:Y\to X$ is quasi-compact and quasi-separated.  Then
\[
f_\tau^*: \Cons_\tau(X,\Lam)\longrightarrow\Cons_\tau(Y,\Lam).
\]
If both $X$ and $Y$ are noetherian, the qcqs hypothesis is automatic.
\dep{Definition 3.5, Proposition 4.2, and preservation of constructible
subsets by inverse image under a spectral, hence qcqs, morphism.}
\end{proposition}

\begin{proposition}[Composition coherence for pullback]
For $Z\xrightarrow{g}Y\xrightarrow{f}X$ there are canonical natural
isomorphisms
\[
c_{g,f}:(f\circ g)_\tau^*\xrightarrow{\sim}g_\tau^*f_\tau^*,
\qquad
u_X:(\id_X)_\tau^*\xrightarrow{\sim}\id.
\]
For three composable morphisms the two composites made from the $c_{-,-}$
agree, and the two unit triangles made from $c_{-,-}$ and $u$ commute.
\dep{Definition 4.1 and associativity and unit coherence of inverse image and
extension of scalars.}
\end{proposition}

\section{Finite etale--pro-etale comparison in the heart}

\begin{definition}[Comparison geometric morphism]
The inclusion of etale test schemes among weakly etale test schemes induces a
geometric morphism
\[
\nu_X:\Sh(X_{\proet})\longrightarrow\Sh(X_{\et}).
\]
Its direct image $\nu_{X,*}$ is restriction to etale objects.  Its inverse
image is denoted $\nu_X^{-1}$.  The canonical coefficient map is an
isomorphism
\[
\nu_X^{-1}\underline\Lam_{X,\et}
\xrightarrow{\sim}\underline\Lam_{X,\proet}.
\]
For module sheaves define
\[
\nu_X^*M:=\underline\Lam_{X,\proet}
\otimes_{\nu_X^{-1}\underline\Lam_{X,\et}}\nu_X^{-1}M,
\quad
\nu_X^*: \mathsf A_\et(X,\Lam)\to\mathsf A_\proet(X,\Lam).
\]
\end{definition}

\begin{lemma}[Fpqc descent for pulled-back etale sections]
Let $F$ be a sheaf of sets on $X_\et$, let $f:Y\to X$, and put
\[
A_F(Y):=\Gamma(Y_\et,f_\et^{-1}F).
\]
For every fpqc covering $\{Y_i\to Y\}$, the sequence
\[
A_F(Y)\longrightarrow\prod_iA_F(Y_i)
\rightrightarrows\prod_{i,j}A_F(Y_i\times_YY_j)
\]
is an equalizer.
\dep{etale pullback from Definition 4.1 and fpqc descent for sections of an
etale sheaf.}
\end{lemma}

\begin{proposition}[Section formula for inverse image]
For every $F\in\Sh(X_\et)$ and every pro-etale object $f:Y\to X$, there is a
functorial identification
\[
(\nu_X^{-1}F)(Y)=\Gamma(Y_\et,f_\et^{-1}F).
\]
For an etale object $U\to X$, this specializes to
$(\nu_X^{-1}F)(U)=F(U)$.
\dep{Definition 5.1 and Lemma 5.2; the displayed sheaf is characterized as the
left adjoint of restriction $\nu_{X,*}$.}
\end{proposition}

\begin{corollary}[Unit isomorphism and full faithfulness]
For every etale sheaf $F$, the adjunction unit
\[
F\longrightarrow\nu_{X,*}\nu_X^{-1}F
\]
is an isomorphism.  Consequently $\nu_X^{-1}$ is fully faithful on sheaves of
sets, abelian groups, and $\Lam$-modules.
\dep{Proposition 5.3 and the adjunction $\nu_X^{-1}\dashv\nu_{X,*}$.}
\end{corollary}

\begin{proposition}[Exactness and coefficient compatibility]
The module functor
\[
\nu_X^*: \mathsf A_\et(X,\Lam)
\longrightarrow\mathsf A_\proet(X,\Lam)
\]
is exact, fully faithful, and strong symmetric monoidal.  It preserves finite
limits, all colimits, kernels, cokernels, and extensions.
\dep{Definition 5.1, Corollary 5.4, and exactness of inverse image for a
geometric morphism.}
\end{proposition}

\begin{theorem}[Torsor comparison]
Let $G$ be a sheaf of groups on $X_\et$.  Pullback induces a bijection
\[
H^1(X_\et,G)\xrightarrow{\sim}
H^1(X_\proet,\nu_X^{-1}G),
\]
where $H^1$ denotes isomorphism classes of torsors.  The assertion remains
valid after replacing $X$ by any etale $X$-scheme.
\dep{Corollary 5.4, the section formula of Proposition 5.3, approximation of
affine pro-etale objects by affine etale objects, and descent of finite cocycle
data.}
\end{theorem}

\begin{theorem}[Finite-lisse essential surjectivity]
Pullback restricts to an equivalence
\[
\left\{\begin{array}{c}
\text{finite locally constant $\Lam$-modules on }X_\et
\end{array}\right\}
\xrightarrow{\ \sim\ }
\left\{\begin{array}{c}
\text{finite locally constant $\Lam$-modules on }X_\proet.
\end{array}\right\}
\]
For a pro-etale local system locally isomorphic to $\underline M$, the descent
object is obtained from the torsor
$\operatorname{Isom}_\Lam(\underline M,F)$ under the finite group
$\operatorname{Aut}_\Lam(M)$.
\dep{Theorem 5.6, Corollary 5.4, and Definition 3.1.}
\end{theorem}

\begin{definition}[Support of a section]
Let $F$ be a sheaf on $U_\tau$, let $W\to U$ be an object of the selected
site, and let $s\in F(W)$.  The zero locus of $s$ is the largest open
subobject of $W$ on which $s$ restricts to zero.  Its complement is the
closed subset $\operatorname{supp}(s)\subseteq |W|$.
\end{definition}

\begin{definition}[Support-closed extension presheaf]
Let $j:U\hookrightarrow X$ be a quasi-compact open immersion and let
$\tau\in\{\et,\proet\}$.  For $F\in\mathsf A_\tau(U,\Lam)$ and an object
$V\to X$, write $V_U=V\times_XU$ and set
\[
 (j^{\mathrm{pre},\tau}_!F)(V)
 =\{s\in F(V_U):\operatorname{supp}(s)\text{ is closed in }V\}.
\]
Pullback of sections defines the restriction maps, and a morphism $F\to G$
defines a morphism
$j^{\mathrm{pre},\tau}_!F\to j^{\mathrm{pre},\tau}_!G$.
\end{definition}

\begin{lemma}[Restriction preserves the support condition]
The preceding restriction maps preserve support-closed sections, identities,
and composition.  Thus $j^{\mathrm{pre},\tau}_!F$ is a presheaf of
$\Lam$-modules on $X_\tau$, functorial in $F$.
\dep{Definitions 5.8 and 5.9 and stability of closed support under pullback
along morphisms in the selected site.}
\end{lemma}

\begin{definition}[Open extension by zero]
Let $a_{X,\tau}$ denote sheafification in $\Lam$-modules on $X_\tau$.  Define
\[
 j_!^\tau F:=a_{X,\tau}(j^{\mathrm{pre},\tau}_!F).
\]
The functoriality of the presheaf construction and of sheafification gives an
additive functor
\[
j_!^\tau:\mathsf A_\tau(U,\Lam)\longrightarrow
\mathsf A_\tau(X,\Lam).
\]
Restriction of support-closed sections gives
$\eta_F:F\to j_\tau^*j_!^\tau F$, and extension of a support-closed section
gives $\epsilon_G:j_!^\tau j_\tau^*G\to G$.
\end{definition}

\begin{definition}[Closed direct image]
Let $i:Z\hookrightarrow X$ be a closed immersion.  Define
\[
i_*^\tau:\mathsf A_\tau(Z,\Lam)\longrightarrow
\mathsf A_\tau(X,\Lam),
\qquad
(i_*^\tau G)(V):=G(V\times_XZ).
\]
\end{definition}

\begin{proposition}[Exactness and adjunctions for an open--closed pair]
Let $j:U\hookrightarrow X$ be open and let $i:Z\hookrightarrow X$ be its
closed complement.  The functors $j_\tau^*$, $j_!^\tau$, $i_\tau^*$, and
$i_*^\tau$ are exact.  One has adjunctions
\[
j_!^\tau\dashv j_\tau^*,\qquad i_\tau^*\dashv i_*^\tau,
\]
the functors $j_!^\tau$ and $i_*^\tau$ are fully faithful, and
\[
j_\tau^*j_!^\tau\simeq\id,\qquad i_\tau^*i_*^\tau\simeq\id,
\qquad i_\tau^*j_!^\tau=0,\qquad j_\tau^*i_*^\tau=0.
\]
On geometric stalks, $j_!^\tau F$ equals $F$ over $U$ and zero over $Z$,
whereas $i_*^\tau G$ equals $G$ over $Z$ and zero over $U$.
\dep{Definitions 5.11 and 5.12 and the open and closed subtopoi associated with
$|X|=|U|\amalg|Z|$.}
\end{proposition}

\begin{proposition}[Open and closed comparison maps]
For an open immersion $j:U\hookrightarrow X$ and an etale module sheaf $G$ on
$U$, the canonical map
\[
\nu_X^*j_!^\et G\xrightarrow{\sim}j_!^\proet\nu_U^*G
\]
is an isomorphism.  For a closed immersion $i:Z\hookrightarrow X$ and an
etale module sheaf $H$ on $Z$, the canonical map
\[
\nu_X^*i_*^\et H\xrightarrow{\sim}i_*^\proet\nu_Z^*H
\]
is an isomorphism.
\dep{Definitions 5.11 and 5.12, Proposition 5.13, Corollary 5.4, adjoint uniqueness for $j_!$, and
the qcqs direct-image comparison for $i_*$.}
\end{proposition}

\begin{proposition}[Constructibility of open and closed extension]
If $j:U\hookrightarrow X$ is a quasi-compact open immersion, then
\[
j_!^\tau:\Cons_\tau(U,\Lam)\longrightarrow\Cons_\tau(X,\Lam).
\]
For every closed immersion $i:Z\hookrightarrow X$,
\[
i_*^\tau:\Cons_\tau(Z,\Lam)\longrightarrow\Cons_\tau(X,\Lam).
\]
Combining an adapted stratification of the source with the complementary
stratum gives an adapted stratification of $X$: the extension is finite locally
constant on every source stratum and zero on the complementary stratum.
\dep{Definition 3.5, Proposition 5.13, and the affine-local two-stratum
decomposition by the image of $i$ and its complement.}
\end{proposition}

\begin{proposition}[Open--closed localization sequence]
Let $j:U\hookrightarrow X$ be open and let $i:Z\hookrightarrow X$ be its
closed complement.  For every $F\in\mathsf A_\tau(X,\Lam)$ there is a
functorial short exact sequence
\[
0\longrightarrow j_!^\tau j_\tau^*F
\longrightarrow F
\longrightarrow i_*^\tau i_\tau^*F
\longrightarrow0.
\]
\dep{Proposition 5.13 and the decomposition of geometric stalks between $U$
and $Z$.}
\end{proposition}

\begin{theorem}[Constructible devissage]
Let $F\in\Cons_\proet(X,\Lam)$.  On every affine open admitting an adapted
finite stratification, the counit
\[
\nu_X^*\nu_{X,*}F\longrightarrow F
\]
is an isomorphism.  Consequently every constructible pro-etale sheaf lies in
the essential image of $\nu_X^*$.
\dep{Theorem 5.7, Propositions 5.14--5.16, Corollary 5.4, and noetherian
induction on the number of strata using a closed stratum and its
quasi-compact open complement.}
\end{theorem}

\begin{theorem}[Equivalence of constructible hearts]
For every scheme $X$ and finite commutative ring $\Lam$, inverse image is an
exact equivalence of abelian categories
\[
\nu_X^*: \Cons_\et(X,\Lam)
\xrightarrow{\ \sim\ }\Cons_\proet(X,\Lam).
\]
A quasi-inverse is the restriction of $\nu_{X,*}$ to constructible objects.
The equivalence preserves finite limits and colimits, tensor products, finite
locally constant objects, and restrictions to locally closed subschemes.
\dep{Proposition 5.5, Theorem 5.17, and Corollary 3.8.}
\end{theorem}

\begin{theorem}[Naturality with respect to scheme pullback]
For every $f:Y\to X$ there is a canonical natural isomorphism
\[
\beta_f:f_\proet^*\nu_X^*\xrightarrow{\sim}
\nu_Y^*f_\et^*
\]
of functors $\mathsf A_\et(X,\Lam)\to\mathsf A_\proet(Y,\Lam)$.
For composable morphisms, the square comparing $\beta_{f\circ g}$ with
$\beta_f$, $\beta_g$, and the composition isomorphisms of Proposition 4.4
commutes; the identity square also commutes.  If the relevant morphisms are
qcqs, these isomorphisms restrict to the constructible hearts.
\dep{Definitions 4.1 and 5.1, Proposition 4.4, and Theorem 5.18 for the
constructible restriction.}
\end{theorem}

\section{Complexes, homotopy, and Verdier localization}

\begin{definition}[Cochain complex and chain homotopy]
Let $\mathcal A$ be an abelian category.  A cochain complex is a family
$K=(K^i,d_K^i)_{i\in\mathbf Z}$ with
$d_K^{i+1}d_K^i=0$.  A chain map $u:K\to L$ is a family $u^i:K^i\to L^i$
such that $d_L^iu^i=u^{i+1}d_K^i$.  Two chain maps $u,v:K\to L$ are homotopic
if there are maps $h^i:K^i\to L^{i-1}$ satisfying
\[
u^i-v^i=d_L^{i-1}h^i+h^{i+1}d_K^i.
\]
Write $\Cpx(\mathcal A)$ for complexes and chain maps and
$\Kcat(\mathcal A)$ for the quotient by chain homotopy.  Superscripts $b$,
$+$, and $-$ denote bounded, bounded-below, and bounded-above complexes,
respectively.
\end{definition}

\begin{definition}[Shift and mapping cone]
The shift is
\[
(K[1])^i=K^{i+1},\qquad d_{K[1]}^i=-d_K^{i+1}.
\]
For a chain map $u:K\to L$, define
\[
\Cone(u)^i=L^i\oplus K^{i+1},\qquad
d_{\Cone(u)}^i=
\begin{pmatrix}d_L^i&u^{i+1}\\0&-d_K^{i+1}\end{pmatrix}.
\]
The triangles in $\Kcat(\mathcal A)$ isomorphic to
\[
K\xrightarrow{u}L\longrightarrow\Cone(u)\longrightarrow K[1]
\]
are the distinguished triangles.
\end{definition}

\begin{definition}[Cohomology, acyclicity, and quasi-isomorphism]
The $i$th cohomology object of $K$ is
\[
H^i(K):=\ker(d_K^i)/\operatorname{im}(d_K^{i-1}).
\]
The complex is acyclic if every $H^i(K)$ is zero.  A chain map $q:K\to L$
is a quasi-isomorphism if every $H^i(q)$ is an isomorphism.  Let $W_\mathcal A$
denote the class of quasi-isomorphisms in $\Kcat(\mathcal A)$.
\end{definition}

\begin{proposition}[Quasi-isomorphism calculus]
The class $W_\mathcal A$ contains identities, is closed under composition,
satisfies two-out-of-three, is invariant under shift, and is saturated.  A map
$q$ lies in $W_\mathcal A$ if and only if $\Cone(q)$ is acyclic.  The acyclic
complexes form a thick triangulated subcategory of $\Kcat(\mathcal A)$.
\dep{Definitions 6.1--6.3 and the long exact cohomology sequence of a cone.}
\end{proposition}

\begin{definition}[Verdier-localized derived categories]
Define
\[
D(\mathcal A):=\Kcat(\mathcal A)[W_\mathcal A^{-1}],\qquad
D^+(\mathcal A):=\Kcat^+(\mathcal A)[(W_\mathcal A\cap\Kcat^+)^{-1}],
\]
\[
D^-(\mathcal A):=\Kcat^-(\mathcal A)[(W_\mathcal A\cap\Kcat^-)^{-1}],
\]
and
\[
D^b(\mathcal A):=\Kcat^b(\mathcal A)[(W_\mathcal A\cap\Kcat^b)^{-1}].
\]
The localization functor $Q$ is universal among functors that invert
quasi-isomorphisms.  Equivalently, this is the Verdier quotient by the thick
subcategory of acyclic complexes.
\end{definition}

\begin{theorem}[Triangulated structure of the localization]
The categories in Definition 6.5 are ordinary triangulated categories.  The
localization functors are exact, shifts are induced from cochain shifts, and
distinguished triangles are precisely the triangles isomorphic to images of
mapping-cone triangles.  For $\mathcal A=\mathsf A_\tau(X,\Lam)$ these derived
categories are locally small.
\dep{Proposition 6.4, Verdier localization, the universe convention, and
Theorem 2.8.}
\end{theorem}

\begin{proposition}[Bounded models]
The canonical functor from the Verdier localization of bounded complexes to
the strictly full subcategory of $D(\mathcal A)$ consisting of objects with
bounded cohomology is an equivalence.  Hence either description may be used
for $D^b(\mathcal A)$.
\dep{Definitions 6.3 and 6.5 and good truncation of complexes.}
\end{proposition}

\begin{definition}[Site-derived categories and cohomology functors]
For $\tau\in\{\et,\proet\}$ put
\[
D_\tau(X,\Lam):=D(\mathsf A_\tau(X,\Lam)),
\quad
D_\tau^\diamond(X,\Lam):=D^\diamond(\mathsf A_\tau(X,\Lam)),
\quad \diamond\in\{+,-,b\}.
\]
Quasi-isomorphism invariance makes
\[
H^i_\tau:D_\tau(X,\Lam)\longrightarrow\mathsf A_\tau(X,\Lam)
\]
well defined by $H^i_\tau(QK)=H^i(K)$.
\end{definition}

\section{The standard t-structure and constructible complexes}

\begin{definition}[Good truncations]
For a complex $K$ and $n\in\mathbf Z$, define
\[
(\tau_{\le n}K)^i=
\begin{cases}
K^i&i<n,\\
\ker(d_K^n)&i=n,\\
0&i>n,
\end{cases}
\qquad
(\tau_{\ge n}K)^i=
\begin{cases}
0&i<n,\\
\operatorname{coker}(d_K^{n-1})&i=n,\\
K^i&i>n.
\end{cases}
\]
The differentials are induced by those of $K$.  These constructions preserve
quasi-isomorphisms and therefore define endofunctors of the derived category.
\end{definition}

\begin{definition}[A t-structure]
A $t$-structure on a triangulated category $\mathcal T$ is a pair of strictly
full subcategories $(\mathcal T^{\le0},\mathcal T^{\ge0})$ such that
\[
\mathcal T^{\le0}[1]\subseteq\mathcal T^{\le0},\qquad
\mathcal T^{\ge0}[-1]\subseteq\mathcal T^{\ge0},\qquad
\Hom(A,B[-1])=0
\]
for $A\in\mathcal T^{\le0}$ and $B\in\mathcal T^{\ge0}$, and every
$K\in\mathcal T$ occurs in a distinguished triangle
\[
A\longrightarrow K\longrightarrow B[-1]\longrightarrow A[1]
\]
with $A\in\mathcal T^{\le0}$ and $B\in\mathcal T^{\ge0}$.  Its heart is
$\mathcal T^{\le0}\cap\mathcal T^{\ge0}$.  It is bounded if every object lies
in $\mathcal T^{\le n}\cap\mathcal T^{\ge m}$ for some $m\le n$.
\end{definition}

\begin{theorem}[Standard bounded t-structure]
Set
\[
D_\tau^{\le0}:=\{K:H^i_\tau(K)=0\text{ for }i>0\},\qquad
D_\tau^{\ge0}:=\{K:H^i_\tau(K)=0\text{ for }i<0\}.
\]
On $D_\tau^b(X,\Lam)$ these subcategories form a bounded $t$-structure.
Its truncation functors are those of Definition 7.1 and its heart is
canonically equivalent to $\mathsf A_\tau(X,\Lam)$ by $M\mapsto M[0]$, with
inverse $H^0_\tau$.
\dep{Definitions 6.8, 7.1, and 7.2 and Theorem 6.6.}
\end{theorem}

\begin{definition}[Constructible derived categories]
Define the strictly full triangulated candidates
\[
D_{c,\tau}(X,\Lam)
:=\{K\in D_\tau(X,\Lam):H^i_\tau(K)\in\Cons_\tau(X,\Lam)\ \forall i\},
\]
\[
D^+_{c,\tau}:=D_{c,\tau}\cap D^+_\tau,
\qquad
 D^-_{c,\tau}:=D_{c,\tau}\cap D^-_\tau,
\qquad
\DbcAt{\tau}:=D_{c,\tau}\cap D^b_\tau.
\]
Thus $K\in\DbcAt{\tau}(X,\Lam)$ means precisely that every cohomology sheaf is
constructible and only finitely many are nonzero.
\end{definition}

\begin{lemma}[Cohomology of a selected triangle]
For every selected distinguished triangle
\[
 K\xrightarrow{u}L\longrightarrow C\longrightarrow K[1]
\]
in $D_\tau(X,\Lam)$ and every $i$, there is a canonical short exact sequence
\[
0\longrightarrow\operatorname{coker}H^i_\tau(u)
\longrightarrow H^i_\tau(C)
\longrightarrow\ker H^{i+1}_\tau(u)\longrightarrow0.
\]
It is natural with respect to morphisms of distinguished triangles.
\dep{the long exact cohomology sequence attached to the cone triangle and
Theorem 7.3.}
\end{lemma}

\begin{proposition}[Triangulated constructible subcategories]
The categories $D_{c,\tau}$, $D^+_{c,\tau}$, $D^-_{c,\tau}$, and
$\DbcAt{\tau}$ are strictly full triangulated subcategories of their ambient
derived categories.  They are
closed under shifts, cones, and direct summands.
\dep{Lemma 7.5 and Theorem 3.7.}
\end{proposition}

\begin{theorem}[Restriction of the standard t-structure]
Good truncations preserve constructible cohomology.  The standard
$t$-structure restricts to $\DbcAt{\tau}(X,\Lam)$ and its heart is
$\Cons_\tau(X,\Lam)$.  The restricted $t$-structure is bounded.
\dep{Definitions 7.1 and 7.4, Theorem 7.3, and Proposition 7.6.}
\end{theorem}

\section{Cohomological and derived site comparison}

\begin{definition}[Derived comparison functors]
The exact functor of Proposition 5.5, applied degreewise to complexes,
preserves quasi-isomorphisms and defines a triangulated functor
\[
\nu_X^*:D_\et(X,\Lam)\longrightarrow D_\proet(X,\Lam),
\]
with restrictions to the bounded-above, bounded-below, and bounded derived
categories.  The same notation is used for the abelian and derived functors.

The restriction functor
\[
\nu_{X,*}:\mathsf A_\proet(X,\Lam)\longrightarrow
\mathsf A_\et(X,\Lam)
\]
is left exact.  For $M\in D^+_\proet(X,\Lam)$ choose a bounded-below
injective resolution $M\to I^\bullet$ in
$\mathsf A_\proet(X,\Lam)$ and define
\[
R\nu_{X,*}M:=\nu_{X,*}I^\bullet,
\qquad
R\nu_{X,*}:D^+_\proet(X,\Lam)\longrightarrow D^+_\et(X,\Lam).
\]
The comparison theorem for injective resolutions makes the result independent
of the resolution and makes $R\nu_{X,*}$ a triangulated functor.  Set
$R^q\nu_{X,*}:=H^q_\et\circ R\nu_{X,*}$ on sheaves placed in degree zero.
\end{definition}

\begin{proposition}[Acyclicity of pulled-back injectives]
Let $I$ be an injective object of $\mathsf A_\et(X,\Lam)$.  Then
\[
R^q\nu_{X,*}(\nu_X^*I)=0\qquad(q>0),
\]
and $\nu_{X,*}\nu_X^*I\simeq I$.  The statement holds after restriction
to every etale object of $X$.
\dep{Definition 8.1, Proposition 5.3, Corollary 5.4, approximation of affine pro-etale objects
by inverse systems of affine etale objects, and exactness of filtered colimits.}
\end{proposition}

\begin{theorem}[Bounded-below derived unit]
For every $K\in D^+_\et(X,\Lam)$, the adjunction unit is an isomorphism
\[
K\xrightarrow{\sim}R\nu_{X,*}\nu_X^*K.
\]
Consequently
\[
\nu_X^*:D^+_\et(X,\Lam)\longrightarrow D^+_\proet(X,\Lam)
\]
is fully faithful.
\dep{Proposition 8.2, bounded-below injective resolutions supplied by Theorem
2.8, and the Leray acyclicity criterion.}
\end{theorem}

\begin{proposition}[Essential image in bounded-below degree]
Let $\mathcal E_X\subseteq\mathsf A_\proet(X,\Lam)$ be the essential image of
$\nu_X^*$.  It is a weak Serre subcategory.  The functor $\nu_X^*$ induces an
equivalence
\[
D^+_\et(X,\Lam)
\xrightarrow{\ \sim\ }
D^+_{\mathcal E_X,\proet}(X,\Lam),
\]
where the target consists of complexes whose cohomology sheaves lie in
$\mathcal E_X$.  A quasi-inverse is $R\nu_{X,*}$.
\dep{Proposition 5.5, Theorem 8.3, and closure of the essential image under
kernels, cokernels, and extensions obtained from the derived unit.}
\end{proposition}

\begin{theorem}[Finite-level derived site comparison]
For every scheme $X$ and finite commutative ring $\Lam$, exact inverse image
induces a $t$-exact triangulated equivalence
\[
\nu_X^*:D^+_{c,\et}(X,\Lam)
\xrightarrow{\ \sim\ }D^+_{c,\proet}(X,\Lam).
\]
It restricts to a $t$-exact triangulated equivalence
\[
\nu_X^*: \DbcAt{\et}(X,\Lam)
\xrightarrow{\ \sim\ }\DbcAt{\proet}(X,\Lam).
\]
For every $i$ there is a natural isomorphism
\[
H^i_\proet(\nu_X^*K)\simeq\nu_X^*H^i_\et(K).
\]
\dep{Theorem 5.18, Proposition 8.4, and Definition 7.4.}
\end{theorem}

\begin{corollary}[Ext comparison]
For $F,G\in\Cons_\et(X,\Lam)$ and every $r\ge0$, pullback induces an
isomorphism
\[
\operatorname{Ext}^r_{X_\et,\Lam}(F,G)
\xrightarrow{\sim}
\operatorname{Ext}^r_{X_\proet,\Lam}(\nu_X^*F,\nu_X^*G),
\]
where both Ext groups are computed in the corresponding ambient module-sheaf
categories.
\dep{Theorem 8.5 and the identification
$\operatorname{Ext}^r(F,G)=\Hom_{D^+}(F,G[r])$.}
\end{corollary}

\section{K-flat tensor products and finite Tor amplitude}

\begin{definition}[Tensor product of complexes]
For complexes $K,L$ of $\underline\Lam_{X,\tau}$-modules, their total tensor
complex has
\[
(K\otimes L)^n=\bigoplus_{p+q=n}K^p\otimes L^q
\]
and differential
\[
d(x\otimes y)=d_Kx\otimes y+(-1)^p x\otimes d_Ly
\qquad(x\in K^p).
\]
\end{definition}

\begin{definition}[K-flat complex]
A complex $P$ is K-flat if $P\otimes A$ is acyclic for every acyclic complex
$A$.  A K-flat replacement of $K$ is a quasi-isomorphism $P\to K$ with $P$
K-flat.
\end{definition}

\begin{theorem}[Existence of K-flat replacements]
Every complex in $\Cpx(\mathsf A_\tau(X,\Lam))$ admits a K-flat replacement.
Replacements may be chosen functorially up to a natural quasi-isomorphism.
Tensoring a K-flat complex with a quasi-isomorphism produces a
quasi-isomorphism.
\dep{Theorem 2.8 and the existence of flat generators for module sheaves on a
small ringed site.}
\end{theorem}

\begin{definition}[Derived tensor product]
Choose K-flat replacements $P(K)\to K$ and $P(L)\to L$.  Define
\[
K\otimes^L_{\underline\Lam_{X,\tau}}L
:=P(K)\otimes_{\underline\Lam_{X,\tau}}P(L).
\]
Theorem 9.3 makes this independent, up to unique derived natural isomorphism,
of the replacements and gives a triangulated bifunctor
\[
\otimes^L:D_\tau(X,\Lam)\times D_\tau(X,\Lam)
\longrightarrow D_\tau(X,\Lam).
\]
\end{definition}

\begin{proposition}[Derived symmetric monoidal coherence]
The ordinary associator, symmetry, and unit maps induce natural isomorphisms
\[
(K\otimes^LL)\otimes^LM\simeq K\otimes^L(L\otimes^LM),
\quad K\otimes^LL\simeq L\otimes^LK,
\quad \underline\Lam_{X,\tau}\otimes^LK\simeq K.
\]
They satisfy the pentagon, hexagon, symmetry, and unit triangle identities.
The functor $K\otimes^L-$ is triangulated for every $K$.
\dep{Definitions 9.1 and 9.4 and Theorem 9.3.}
\end{proposition}

\begin{definition}[Finite Tor amplitude]
An object $P\in D_\tau(X,\Lam)$ has Tor amplitude in $[a,b]$ if, for every
module sheaf $M$ placed in degree zero,
\[
H^i_\tau(P\otimes^L M)=0\qquad(i\notin[a,b]).
\]
It has finite Tor amplitude if such integers $a\le b$ exist.  Define
\[
\DctfAt{\tau}(X,\Lam)
:=\{P\in\DbcAt{\tau}(X,\Lam):P\text{ has finite Tor amplitude}\}.
\]
\end{definition}

\begin{definition}[Stratified perfect complex]
An object $P\in\DbcAt{\tau}(X,\Lam)$ is stratified perfect if, affine locally on
$X$, one finite constructible stratification makes every restriction of $P$
quasi-isomorphic to a bounded complex of finite projective $\Lam$-modules
viewed as constant sheaves.
\end{definition}

\begin{proposition}[Constructibility of Tor sheaves]
For $F,G\in\Cons_\tau(X,\Lam)$ and $r\ge0$, the sheaf
\[
\mathcal T\!or_r^{\underline\Lam_{X,\tau}}(F,G)
\]
is constructible.  On a common stratum where $F=\underline M$ and
$G=\underline N$, it is the constant sheaf associated with
$\Tor_r^\Lam(M,N)$.
\dep{Lemma 3.2, Proposition 3.6, Lemma 1.3, and finite free resolutions degree
by degree.}
\end{proposition}

\begin{theorem}[Bounded-above constructibility of derived tensor]
If $K,L\in D^-_{c,\tau}(X,\Lam)$, then
\[
K\otimes^L_{\underline\Lam_{X,\tau}}L
\in D^-_{c,\tau}(X,\Lam).
\]
Equivalently, every cohomology sheaf of the tensor is constructible and the
cohomology vanishes in sufficiently high degrees.
\dep{Proposition 9.8, the finite hyper-Tor filtrations obtained from good
truncations, and extension closure in Theorem 3.7.}
\end{theorem}

\begin{corollary}[Bounded tensor with a finite-Tor factor]
If $K,L\in\DbcAt{\tau}(X,\Lam)$ and one factor has finite Tor amplitude, then
\[
K\otimes^LL\in\DbcAt{\tau}(X,\Lam).
\]
If $K$ has Tor amplitude $[a,b]$ and $H^i(L)=0$ outside $[r,s]$, then the tensor
has cohomology only in $[a+r,b+s]$.  In particular, tensoring with a
stratified perfect complex preserves $\DbcAt{\tau}$.
\dep{Theorem 9.9 and Definitions 9.6 and 9.7.}
\end{corollary}

\begin{proposition}[No unrestricted bounded tensor closure]
For a general finite torsion ring, $\DbcAt{\tau}(X,\Lam)$ is not closed under
derived tensor.  For example, with $X$ a geometric point,
$\Lam=\mathbf Z/\ell^2\mathbf Z$, and $M=\Lam/\ell\Lam$, one has
\[
\Tor_i^\Lam(M,M)\simeq M\qquad(i\ge0),
\]
so $M\otimes^L_\Lam M$ is unbounded although both inputs lie in degree zero.
\dep{Definition 9.4 and the periodic free resolution of $M$ with differential
multiplication by $\ell$.}
\end{proposition}

\section{Derived pullback and comparison coherences}

\begin{definition}[Derived same-coefficient pullback]
Termwise pullback preserves quasi-isomorphisms.  Define
\[
Lf_\tau^*:D_\tau(X,\Lam)\longrightarrow D_\tau(Y,\Lam)
\]
by applying $f_\tau^*$ to a representative complex.  There is a canonical
identification $Lf_\tau^*=f_\tau^*$ on the derived category.
\end{definition}

\begin{proposition}[Cohomology and constructibility under derived pullback]
For every $f:Y\to X$ and every $K\in D_\tau(X,\Lam)$,
\[
H^i_\tau(Lf_\tau^*K)\simeq f_\tau^*H^i_\tau(K).
\]
The functor $Lf_\tau^*$ is $t$-exact.  If $f$ is qcqs, it restricts to
\[
Lf_\tau^*:D^\diamond_{c,\tau}(X,\Lam)
\longrightarrow D^\diamond_{c,\tau}(Y,\Lam),
\qquad \diamond\in\{+,-,b\}.
\]
\dep{Propositions 4.2 and 4.3, Definition 10.1, and Theorem 7.3.}
\end{proposition}

\begin{proposition}[Derived pullback composition coherence]
The underived isomorphisms of Proposition 4.4 descend to
\[
L(f\circ g)_\tau^*\xrightarrow{\sim}Lg_\tau^*Lf_\tau^*,
\qquad L(\id_X)_\tau^*\xrightarrow{\sim}\id.
\]
They satisfy the same associativity pentagon and unit triangles as the
underived pullbacks.
\dep{Proposition 4.4 and Definition 10.1.}
\end{proposition}

\begin{theorem}[Pullback--tensor compatibility]
For every $f:Y\to X$ and $K,L\in D_\tau(X,\Lam)$ there is a natural
isomorphism
\[
Lf_\tau^*(K\otimes^L L)
\xrightarrow{\sim}
Lf_\tau^*K\otimes^L Lf_\tau^*L.
\]
It is compatible with the associator, symmetry, and unit of Proposition 9.5
and with the composition maps of Proposition 10.3.  The isomorphism restricts
to $\Dbc$ whenever both sides belong to the bounded constructible category,
in particular in the finite-Tor situation of Corollary 9.10.
\dep{Proposition 4.2, Theorem 9.3, Proposition 9.5, and Definition 10.1.}
\end{theorem}

\begin{theorem}[Derived etale--pro-etale pullback square]
For every morphism of schemes $f:Y\to X$, the natural isomorphism of Theorem
5.16 descends to an isomorphism of triangulated functors
\[
Lf_\proet^*\nu_X^*
\xrightarrow{\sim}
\nu_Y^*Lf_\et^*:
D_\et(X,\Lam)\longrightarrow D_\proet(Y,\Lam).
\]
It is compatible with identity morphisms and composition of scheme maps.  For
qcqs $f$ it restricts to the constructible bounded and bounded-below
subcategories.
\dep{Theorem 5.19, Definition 10.1, Proposition 10.3, and Theorem 8.5.}
\end{theorem}

\begin{theorem}[Derived etale--pro-etale tensor square]
For $K,L\in D_\et(X,\Lam)$ there is a natural isomorphism
\[
\nu_X^*(K\otimes^L_{\underline\Lam_{X,\et}}L)
\xrightarrow{\sim}
\nu_X^*K\otimes^L_{\underline\Lam_{X,\proet}}\nu_X^*L.
\]
It respects the symmetric monoidal coherence maps and is compatible with the
scheme-pullback square of Theorem 10.5.  It restricts to $\Dbc$ in the
finite-Tor range of Corollary 9.10.
\dep{Proposition 5.5, Theorem 9.3, Proposition 9.5, and Theorem 10.4.}
\end{theorem}

\begin{definition}[Derived change of finite coefficients]
For a homomorphism of finite commutative rings $\varphi:\Lam\to\Lam'$, define
\[
\operatorname{ext}_{\varphi,\tau}(K)
:=K\otimes^L_{\underline\Lam_{X,\tau}}
\underline{\Lam'}_{X,\tau},
\qquad
D_\tau(X,\Lam)\to D_\tau(X,\Lam').
\]
No boundedness assertion is included in this definition.
\end{definition}

\begin{proposition}[Coefficient-change comparison and coherence]
For every $\varphi:\Lam\to\Lam'$ there is a natural isomorphism
\[
\nu_X^*\operatorname{ext}_{\varphi,\et}
\xrightarrow{\sim}
\operatorname{ext}_{\varphi,\proet}\nu_X^*.
\]
For $\Lam\to\Lam'\to\Lam''$, these isomorphisms commute with the associativity
isomorphism for iterated derived extension of scalars and with identity ring
maps.  They also commute with the scheme-pullback square of Theorem 10.5.
The functors restrict to $\Dbc$ only under a finite-Tor hypothesis ensuring
boundedness; finite rings alone do not provide that hypothesis.
\dep{Definition 10.7, Theorems 10.5 and 10.6, and the associativity coherence
of Proposition 9.5.}
\end{proposition}

\section{Finite-torsion foundation theorem}

\begin{theorem}[Finite-torsion pro-etale foundation]
Let $X$ be a scheme and let $\Lam$ be a finite commutative ring.  Then:
\begin{enumerate}
\item the categories $\DbcAt{\et}(X,\Lam)$ and
$\DbcAt{\proet}(X,\Lam)$ are explicitly defined Verdier localizations of
ordinary homotopy categories, cut out by bounded constructible cohomology;
\item each carries the bounded standard $t$-structure with heart
$\Cons_\tau(X,\Lam)$;
\item exact inverse image gives a $t$-exact triangulated equivalence
\[
\nu_X^*: \DbcAt{\et}(X,\Lam)\xrightarrow{\sim}
\DbcAt{\proet}(X,\Lam);
\]
\item same-coefficient pullback is an exact triangulated functor and preserves
bounded constructibility for qcqs scheme maps;
\item the ambient derived tensor product is defined by K-flat replacements,
and it restricts to bounded constructible objects whenever one factor has
finite Tor amplitude;
\item all comparison, pullback, tensor, identity, composition, and finite
coefficient-change isomorphisms satisfy the stated coherence laws.
\end{enumerate}
\dep{Theorems 6.6, 7.7, 8.5, 9.3, 10.4--10.6, and Proposition 10.8.}
\end{theorem}

\section{References}

The site and comparison statements are aligned with the Stacks Project,
\href{https://stacks.math.columbia.edu/tag/0988}{Section 61.12},
\href{https://stacks.math.columbia.edu/tag/099R}{Section 61.19}, and
\href{https://stacks.math.columbia.edu/tag/09AI}{Section 61.27}.  In particular,
the section formula is Tag 0GLZ, finite-lisse comparison is Tag 099Y,
bounded-below full faithfulness is contained in Section 61.19, and the
constructible derived equivalence is Lemma 61.27.4.  Verdier localization,
good truncation, K-flat resolutions, and derived tensor use the standard
ordinary derived-category formalism for Grothendieck abelian categories.

\end{document}
